\documentclass[12pt]{article}
\usepackage[a3paper, margin=1in]{geometry}
\usepackage{amsmath, amssymb, amsthm, ulem, booktabs}

\begin{document}

\section*{Domácí úkol 9.1}

Vypracoval: Daniel \textit{``Randál"} Ransdorf \hfill Podpis: \rule{4cm}{0.4pt}

\begin{enumerate}
  
  \item Nejprve určíme dimenzi $Im\,A_{a,b}$, což už víme, že je hodnost matice. V závislosti na parametrech, poté dle teorie určíme zbytek.
    Hledejme tedy nejprve hodnost matice $A_{a,b}$. Upravujme matici elementárními řádkovými a sloupcovými úpravami, o kterých víme, že nemění hodnost matice:

    \[
      \left( \begin{array}{cccc}
        a+2 & b & 1 & 2 \\
        2a+1 & 2b & a+3 & 1 \\
        3a+b+3 & 3b & a+4 & b+3
      \end{array} \right)
      \overset{\text{1.s } + (4\cdot)\text{4.s}}{\sim}
      \left( \begin{array}{cccc}
        a & b & 1 & 2 \\
        2a & 2b & a+3 & 1 \\
        3a & 3b & a+4 & b+3
      \end{array} \right)
    \]

    Vidím, že první a druhý sloupec jsou na sobě závislé. Nemůžeme ale jakýkoli škrtnout, protože bychom mohli ztratit
    údaje, kdyby ten druhý neškrtnutý byl nulový. Vyšetřeme proto tři případy $a\neq0, b\neq0, a=b=0$:
    \begin{enumerate}
      \item $a\neq0$: Když je $a$ nenulové, pak má inverz, a tedy druhý sloupec lze napsat jako násobek prvního.
        Vynulujme ten druhý sloupec přičtením $-a^{-1}b$ násobku prvního sloupce.

        \[
          \sim
          \left( \begin{array}{cccc}
            a & 0 & 1 & 2 \\
            2a & 0 & a+3 & 1 \\
            3a & 0 & a+4 & b+3
          \end{array} \right)
          \overset{\text{1.s } \cdot a^{-1}}{\sim}
          \left( \begin{array}{cccc}
            1 & 0 & 1 & 2 \\
            2 & 0 & a+3 & 1 \\
            3 & 0 & a+4 & b+3
          \end{array} \right)
          \overset{\text{3.s } - \text{1.s}}{\sim}
          \left( \begin{array}{cccc}
            1 & 0 & 0 & 2 \\
            2 & 0 & a+1 & 1 \\
            3 & 0 & a+1 & b+3
          \end{array} \right)
          \overset{\text{4.s } + (3\cdot) \text{1.s}}{\sim}
          \left( \begin{array}{cccc}
            1 & 0 & 0 & 0 \\
            2 & 0 & a+1 & 2 \\
            3 & 0 & a+1 & b+2
          \end{array} \right)
        \]
        \[
          \overset{\text{3.ř } -\text{2.ř}}{\sim}
          \left( \begin{array}{cccc}
            1 & 0 & 0 & 0 \\
            2 & 0 & a+1 & 2 \\
            1 & 0 & 0 & b
          \end{array} \right)
          \overset{\text{2.ř } +(3\cdot)\text{1.ř}}{\sim}
          \left( \begin{array}{cccc}
            1 & 0 & 0 & 0 \\
            0 & 0 & a+1 & 2 \\
            1 & 0 & 0 & b
          \end{array} \right)
          \overset{\text{3.ř } +(4\cdot)\text{1.ř}}{\sim}
          \left( \begin{array}{cccc}
            1 & 0 & 0 & 0 \\
            0 & 0 & a+1 & 2 \\
            0 & 0 & 0 & b
          \end{array} \right)
        \]
        
        Kdybychom na tuto matici chtěli použít Gaussovu eliminaci, tak první řádek už je finální,
        ani nám jakkoli nepomůže (má pivot a zbytek hodnot 0).
        Rozdělme tedy všechny případy:
        \begin{enumerate}
          \item $b=0$: Matice je odstupňovaná a má dva nenulové řádky, hodnost je tedy \underline{2 pro $a\neq0, b=0$}
            \[
              \left( \begin{array}{cccc}
                1 & 0 & 0 & 0 \\
                0 & 0 & a+1 & 2 \\
                0 & 0 & 0 & 0
              \end{array} \right)
            \]
          \item $a+1=0$: Pak třetí řádek zeliminujeme přičtením $-2^{-1}b$ násobku druhého. Tím pádem je hodnost matice
            \underline{2 pro $a\neq0,a=4$}
            \[
              \left( \begin{array}{cccc}
                1 & 0 & 0 & 0 \\
                0 & 0 & 0 & 2 \\
                0 & 0 & 0 & b
              \end{array} \right) \sim \left( \begin{array}{cccc}
                1 & 0 & 0 & 0 \\
                0 & 0 & 0 & 2 \\
                0 & 0 & 0 & 0
              \end{array} \right)
            \]
          \item $b\neq0,a+1\neq0$: Pak $b$ i $a+1$ mají inverz, tedy můžeme pivot $b$ převést na $1$ a ve druhém řádku můžeme vynulovat čtvrtý sloupec pomocí
            přičtením $3$ násobku třetího řádku. Pak vynásobením druhého řádku inverzem $a+1$ získáme odstupňovanou Gauss--Jordanovsky eliminovanou matici se třemi nenulovými řádky,
            hodnost je tedy \underline{3 pro $a\neq0, a\neq4, b\neq0$}.

            \[
              \left( \begin{array}{cccc}
                1 & 0 & 0 & 0 \\
                0 & 0 & a+1 & 2 \\
                0 & 0 & 0 & b
              \end{array} \right)
              \sim
              \left( \begin{array}{cccc}
                1 & 0 & 0 & 0 \\
                0 & 0 & a+1 & 2 \\
                0 & 0 & 0 & 1
              \end{array} \right)
              \sim
              \left( \begin{array}{cccc}
                1 & 0 & 0 & 0 \\
                0 & 0 & a+1 & 0 \\
                0 & 0 & 0 & 1
              \end{array} \right)
              \sim
              \left( \begin{array}{cccc}
                1 & 0 & 0 & 0 \\
                0 & 0 & 1 & 0 \\
                0 & 0 & 0 & 1
              \end{array} \right)
            \]
        \end{enumerate}
      
      \item $b\neq0$: Když je $b$ nenulové, pak má inverz, a tedy první sloupec lze napsat jako násobek druhého.
        Vynulujme ten první sloupec přičtením $-b^{-1}a$ násobku druhého sloupce.

        \[
          \left( \begin{array}{cccc}
            a & b & 1 & 2 \\
            2a & 2b & a+3 & 1 \\
            3a & 3b & a+4 & b+3
          \end{array} \right)
          \sim
          \left( \begin{array}{cccc}
            0 & b & 1 & 2 \\
            0 & 2b & a+3 & 1 \\
            0 & 3b & a+4 & b+3
          \end{array} \right)
          \overset{\text{2.s } \cdot b^{-1}}{\sim}
          \left( \begin{array}{cccc}
            0 & 1 & 1 & 2 \\
            0 & 2 & a+3 & 1 \\
            0 & 3 & a+4 & b+3
          \end{array} \right)
          \overset{\text{1.s } \leftrightarrow \text{ 2.s}}{\sim}
          \left( \begin{array}{cccc}
            1 & 0 & 1 & 2 \\
            2 & 0 & a+3 & 1 \\
            3 & 0 & a+4 & b+3
          \end{array} \right)
        \]

        Tento tvar jsme už jednou viděli a víme, že ho můžeme elementárními úpravami převést do následujícího tvaru
        nezávisle na $a$ (nepoužili jsme nikde $a^{-1}$):

        \[
          \left( \begin{array}{cccc}
            1 & 0 & 0 & 0 \\
            0 & 0 & a+1 & 2 \\
            0 & 0 & 0 & b
          \end{array} \right)
        \]

        $b$ je nenulové, pak tedy můžeme prvek čtvrtého sloupce druhého řádku vynulovat jako v 1.a.iii.:

        \[
          \left( \begin{array}{cccc}
            1 & 0 & 0 & 0 \\
            0 & 0 & a+1 & 0 \\
            0 & 0 & 0 & 1
          \end{array} \right)
        \]

        Pak je zjevně hodnost \underline{2 pro $b\neq0, a=4$}, \underline{3 pro $b\neq0, a\neq4$} -- což nám zesiluje tvrzení z 1.a.iii.,
        protože pro hodnost 3 je přípustné $a=0$.
      
      \item $a=b=0$: Dosaďme nuly a zeliminujme matici:
          \[
            \left( \begin{array}{cccc}
              a & b & 1 & 2 \\
              2a & 2b & a+3 & 1 \\
              3a & 3b & a+4 & b+3
            \end{array} \right)
            \sim
            \left( \begin{array}{cccc}
              0 & 0 & 1 & 2 \\
              0 & 0 & 3 & 1 \\
              0 & 0 & 4 & 3
            \end{array} \right)
            \sim 
            \left( \begin{array}{cccc}
              0 & 0 & 1 & 2 \\
              0 & 0 & 0 & 0 \\
              0 & 0 & 0 & 0
            \end{array} \right)
          \]

        Hodnost je tedy \underline{1 pro $a=b=0$}.
 
    \end{enumerate}
 
        Dejme k sobě podmínky pro jednotlivé hodnoty hodnosti a zjednodušme je:

        \[Rank(A_{a,b}) = \begin{cases}
          1: \quad &\underline{a=0 \land b=0 }\\
          2: \quad &(a\neq0 \land b=0) \;\lor\; (a\neq0 \land a=4) \;\lor\; (b\neq0 \land a=4) \\
                  &\Leftrightarrow \; (a\neq0 \land b=0) \;\lor\; a=4 \;\lor\; (b\neq0 \land a=4) \\
                  &\Leftrightarrow \; \underline{(a\neq0 \land b=0) \;\lor\; a=4 }\\
          3: \quad &(a\neq0 \land a\neq4 \land b\neq0) \;\lor\; (b\neq0 \land a\neq4) \\
                  &\Leftrightarrow \; \underline{b\neq0 \land a\neq4}
        \end{cases}\]

        Z teorie víme, $Rank(A)=Rank(A^T)$. Jelikož $dim(Im(A))=Rank(A)$, pak $dim(Im(A^T))=Rank(A^T)=Rank(A)=dim(Im(A))$.
        Z toho plyne:
        \[dim(Im(A_{a,b}^T)) = dim(Im(A_{a,b})).\]
        
        Dále víme, že $dim(Im(A))$ je rovna počtu bázových sloupců (lineárně nezávislých).
        Také víme, že $dim(Ker(A))$ je rovna počtu volných proměnných (nebázových sloupců -- lineárně závislých).
        V matici $A_{a,b}$ máme 4 sloupce, tedy platí $dim(Im(A_{a,b})) + dim(Ker(A_{a,b})) = 4$.
        V matici $A_{a,b}^T$ jsou 3 sloupce, stejnou úvahou dostaneme $dim(Im(A_{a,b}^T)) + dim(Ker(A_{a,b}^T)) = 3$.
        Odečtením dimenzí obrazu dostáváme:

        \begin{align*}
           dim(Ker(A_{a,b})) = 4 - dim(Im(A_{a,b})) \\
           dim(Ker(A_{a,b}^T)) = 3 - dim(Im(A_{a,b}^T))
        \end{align*}

        Ted' jednoduchými počty tyto hodnoty získáme a vytvoříme výslednou tabulku:


        \begin{table}[h]
          \centering
          \begin{tabular}{lccc}
            \toprule
            & $a=0 \land b=0$ & \;\;$(a\neq0 \land b=0) \lor a=4$\;\; & $b\neq0 \land a\neq4$ \\ 
            \midrule
            $dim(Im(A_{a,b}))$ & 1 & 2 & 3 \\
            $dim(Im(A_{a,b}^T))$ & 1 & 2 & 3 \\
            $dim(Ker(A_{a,b}))$ & 3 & 2 & 1 \\
            $dim(Ker(A_{a,b}^T))$ & 2 & 1 & 0 \\
            \bottomrule
          \end{tabular}
        \end{table}
      
      Ještě na závěr pro jistotu ověřme, že každá dvojice $a,b\in\mathbb{Z}_5$ by měla spadnout do právě jedné z kategorií.
      Kategorie jsou na první pohled bez průniku, myšleno každá dvojice čísel zapadá nejvýše do jedné kategorie. 
      Ještě zbývá ukázat, že jsme žádné dvojice čísel nevynechali --- každá dvojice zapadá minimálně do jedné kategorie.
      Tato výroková formule by měla být pravdivá:
      
      \[
        \begin{aligned}
        & (a = 0 \land b = 0)
          \;\lor\;
          ((a \neq 0 \land b = 0) \lor a = 4)
          \;\lor\;
          (b \neq 0 \land a \neq 4)
        \\
        \Longleftrightarrow\quad &
          (a = 0 \land b = 0)
          \;\lor\;
          (a \neq 0 \land b = 0)
          \;\lor\;
          (a = 4)
          \;\lor\;
          (b \neq 0 \land a \neq 4)
        \\
        \Longleftrightarrow\quad &
          \left((a = 0 \land b = 0)
          \;\lor\;
          (a \neq 0 \land b = 0)\right)
          \;\lor\;
          (a = 4)
          \;\lor\;
          (b \neq 0 \land a \neq 4)
        \\
        \Longleftrightarrow\quad &
          \bigl(\underbrace{(a = 0 \lor a \neq 0)}_{\text{vždy pravda}} \;\land\; b = 0\bigr)
          \;\lor\;
          (a = 4)
          \;\lor\;
          (b \neq 0 \land a \neq 4)
        \\
        \Longleftrightarrow\quad &
          (b = 0)
          \;\lor\;
          (a = 4)
        \;\lor\;
        (b \neq 0 \land a \neq 4)
        \\
        \Longleftrightarrow\quad &
          (b = 0)
          \;\lor\;
          \left((a = 4) \lor (b \neq 0 \land a \neq 4)\right)
        \\
        \Longleftrightarrow\quad &
          (b = 0)
          \;\lor\;
          \left((a = 4) \lor b \neq 0\right)
        \\
        \Longleftrightarrow\quad &
          \underbrace{(b = 0)
          \;\lor\;
          (b \neq 0)}_{\text{vždy pravda}}
          \;\lor\;
          (a = 4)
        \\
        \Longleftrightarrow\quad &
        \text{pravda} \quad \checkmark
        \end{aligned}
      \]



\end{enumerate}

\end{document}
